\section{Design decisions: coordination}
\label{sec:decisions-coordination}

\subsection{Why Hazelcast}

The services need four things from their coordination: queues of work that several replicas can share, a state per
book that several processes change at the same time, one lock (the single index writer) and one event (a new version).
Table~\ref{tab:alternatives} compares the options we considered.

\begin{table}[H]
\centering\small
\begin{tabularx}{\linewidth}{@{}L{3.4cm}X@{}}
\toprule
\textbf{Alternative} & \textbf{Why we did not choose it} \\
\midrule
Control files of Stage 1 & Plain text files do not survive several processes writing at the same time \\
Orchestrator with SQLite & It works, but the orchestrator calls each replica and must know its address; with shared
  queues, adding a replica changes nothing in the Control \\
Broker (RabbitMQ, Kafka) & Acknowledgements and persistence, but one more technology, and no shared state, locks or
  compare-and-set: the state of each book would need yet another store \\
MongoDB for coordination too & It could keep the state of each book, but it has no blocking queues, topics or
  distributed locks; polling a collection for work would load the database that Search reads \\
Shared memory (\code{mmap}) & Only exists on one machine: the assignment would force a redesign for Stage 3 \\
\textbf{Hazelcast} & Queues, a map with compare-and-set and queries, transactions, locks and a reliable topic in one
  dependency. It is used the same way on one machine and on a cluster: for Stage 3 the configuration changes, not the code \\
\bottomrule
\end{tabularx}
\caption{Coordination alternatives.}
\label{tab:alternatives}
\end{table}

Hazelcast has three limits, which shaped the rest of the design:
\begin{enumerate}
  \item \textbf{It is only in memory.} If both members stop, every state is lost. Hence the principle that a state is
    written only after its evidence is on disk, and the recovery of the Control (Section~\ref{sec:recovery}).
  \item \textbf{Its queues have no acknowledgements.} An id taken by a worker that dies is gone from the queue. Hence the
    leases in \code{book-state}, reviewed by the Control (Section~\ref{sec:ownership}).
  \item \textbf{Its strongly consistent primitives are not free.} The CP Subsystem, and with it \code{FencedLock}, is an
    Enterprise feature in 5.7: \code{getCPSubsystem().getLock()} throws \code{UnsupportedOperationException} (measured
    before writing the contract). The index writer uses \code{IMap.lock} instead, which has no fencing token.
\end{enumerate}

\subsection{How the shared state is protected}

\textbf{One library.} Every access to Hazelcast is a method of \code{Coordination} in \code{common}, one per
transition of the contract. No service names a structure or builds a value by hand, so the rules below hold
everywhere by construction, and they are tested once against embedded members.

\textbf{Compare-and-set, not locks.} A transition reads the entry, builds the next value and writes it with
\code{replace(id, expected, next)}, which the member applies only if the stored bytes are still the expected ones.
Two workers that race for a book cannot both win, and nobody holds a lock while it downloads. We did not use an
\code{EntryProcessor}, which would run our code on the members: the members run the official image without our
classes, and the values are Compact records that the members can store, compare and query field by field without them.

\textbf{Transactions where a queue is involved.} Claiming a book (poll and \code{INGESTING}) and completing an
ingestion (\code{DOWNLOADED} and the offer to \code{index-queue}) each touch a map and a queue. They run in one
\code{TWO\_PHASE} transaction, which survives the death of a member during the commit. Without it, a crash between the
two operations would leave a book \code{DOWNLOADED} that no queue holds, or an id taken from the queue that no state
records.

\textbf{Duplicates are harmless.} The queues are work to do, \code{book-state} is the truth. An id can be twice in a
queue (a requeue racing with a late completion, a seed repeated by hand): whoever takes it checks its state in the
claim and drops it if it is not waiting. Together with idempotent work (atomic files, upserts, \code{\$addToSet},
\code{putIfAbsent}), this is why at-least-once delivery is enough and exactly-once is never needed.

\textbf{Fencing where it matters.} \code{IMap.lock} has no fencing token: a writer that froze and comes back late is not
stopped by Hazelcast. For the index this costs at most repeated work, because its writes are idempotent and its claims
fail once the Control has requeued its books. A publication is not idempotent (a late writer could move \code{CURRENT}
back), so Indexing uses the version itself as the token: \code{CURRENT} only moves from V to V+1 with a conditional
update, and a publication whose update matches nothing is aborted.

\subsection{Single machine, replicas, threads and ports}

\textbf{Two Hazelcast members on one host.} One member would be enough on one machine; two show the behaviour of the
cluster that Stage 3 needs: partitions spread over members, one backup of every entry, a member that dies without losing
anything, and clients that move to the other member. On one host they share its failures, so this is evidence of the
mechanism, not real high availability.

\textbf{Ingestion: replicas and workers.} Ingestion scales by adding processes (\code{-{}-scale ingestion=N}) or threads
(\code{INGESTION\_WORKERS}). Both pull from the same queue, so the work is shared without any configuration, and
correctness does not depend on the counts. Downloading is mostly waiting on the network, which favours several workers
per process (4 by default); the limit is Gutenberg, which throttles a client that downloads too fast. With the local
source the work is CPU and disk, and each container has one CPU, so beyond a few workers only more replicas help.

\textbf{One Control.} It is not in the path of any book, so one is enough and its absence only delays leases and
publications. A second one would be safe, because every transition is a compare-and-set and the publication is
idempotent, but it would double the lease queries and the calls to Indexing for nothing. Stage 3 needs a leader election
instead (Section~\ref{sec:stage3}).

\textbf{Ports.} Only nginx is published to the network. The Control, the Hazelcast members and MongoDB have no
authentication, so they are published on \code{127.0.0.1} only, for the administrator, the tests and the services run
from an IDE. Ingestion and Indexing are reached only inside the Docker network. The clients use \code{SINGLE\_MEMBER}
routing: the members announce their Docker-internal addresses, which a process on the host cannot reach, so a client
talks only to the addresses it was given and moves to the next one if its member dies.

\subsection{Where replication helps and where it contends}

\begin{itemize}
  \item \textbf{\code{book-state} is partitioned} by book id over the 271 partitions of the cluster, so claims of
    different books touch different partitions and do not contend.
  \item \textbf{Each queue lives in one partition}, on one member: every claim of every worker goes through it. This
    serializes the claims, and it is the first point of contention of the coordination when the work per book is small
    (the local source). With Gutenberg the download is far longer than a claim, so it should not show there.
  \item \textbf{Each book costs two transactions} (claim and completion), each with the round trips of a two-phase
    commit. That is the coordination cost per book, measured on its own with JMH and inside the running system with the
    metrics of CF-12.
  \item \textbf{The datalake is one disk.} More Ingestion replicas write to the same file system; each book is three
    small files and three atomic moves in the same range folder.
  \item \textbf{Indexing has a single writer by design.} More Indexing replicas would only be standbys; its parallelism is
    in the threads that parse, inside one process.
\end{itemize}

\textbf{Measured.} With JMH (Section~\ref{sec:benchmarks-ingestion}), claiming and completing a book costs 1.6~ms with
one worker, 3.4 times the local ingestion of the book, and 0.36~ms per book with 8 workers, because concurrent workers
overlap their round trips. So with small books the coordination, not the work, sets the rate of a replica; with the
download from Project Gutenberg it is the other way round. The contention of several replicas on the single partition of
each queue was not measured: the integration experiments of Ingestion were not run.

\subsection{Failure tests}
\label{sec:failure-tests}

The scenarios of \code{scripts/chaos/} (CF-11) cause each failure in the middle of a real run on Docker Compose and check
the end state against the contract. Each one leaves a \code{report.md} with what was expected, what happened and the
evidence (the states, a timeline of the books and the logs of every service).

\begin{table}[H]
\centering\small
\begin{tabularx}{\linewidth}{@{}L{3.3cm}XL{3.4cm}@{}}
\toprule
\textbf{Failure} & \textbf{Expected} & \textbf{Result} \\
\midrule
SIGKILL of one of two Ingestion replicas while it ingests & Its books are requeued when their lease expires and the
  other replica finishes them; no book \code{FAILED} or duplicated & \textbf{PASS.} Killed while holding a book; its lease expired and the Control requeued it; 2,000 of 2,000 \code{DOWNLOADED} \\
SIGKILL of one Hazelcast member & Nothing lost (backup), the clients carry on, no recovery; the member rejoins &
  \textbf{PASS.} \code{hz-2} carried on alone with 1,848 books to go; no rebuild; \code{hz-1} rejoined \\
SIGKILL of the Control while Indexing indexes & The workers keep going, nothing is published until it is back; it rebuilds
  nothing & \textbf{PASS.} Indexing went from 100 to 500 books, \code{index\_books} stayed at 100; no rebuild \\
Both members killed and started empty & The Control rebuilds \code{book-state} from the evidence; the books without
  evidence are sent again; none with \code{\_READY} is ingested again; all end \code{PUBLISHED} &
  \textbf{PASS.} Rebuilt in 1.9~s: 3 \code{PUBLISHED}, 100 \code{INDEXED}, 336 \code{DOWNLOADED}; the 1,561 without evidence sent again, none ingested twice \\
\bottomrule
\end{tabularx}
\caption{Documented failure tests. Each run checks that every book reaches its final state exactly once, that none is
\code{FAILED}, and that every book is complete in the datalake and, with Indexing, in \code{index\_log} and
\code{index\_books}.}
\label{tab:failure-tests}
\end{table}

\textbf{How they ran.} On the reference machine, on 8 November, with 2,000 synthetic books per scenario, leases of
15~s and the commit recorded in each \code{report.md} (\code{scripts/chaos/results/}). With the local source a book takes
milliseconds, so the scenarios that kill during the ingestion run one worker per replica and check that books were still
in flight right after the failure. To be sure that the killed Ingestion replica held a book, it is frozen first
(\code{docker pause}) and killed only while the Control shows its book \code{INGESTING}. A first run gave three FAILs
with every check of the end state passing: two came from errors in the checks themselves, and in the third the killed
replica held no book, so no lease was tested. The script was fixed and the four scenarios run again; the table gives
that second run.
